% ════════════════════════════════════════════════════════════════════════════
\section{The Deadly-Race impossibility chain}
\label{sec:chain}
% ════════════════════════════════════════════════════════════════════════════

\begin{flushright}
\itshape ``Nothing was your own except the few cubic centimetres inside your
skull.''\\
\upshape\footnotesize --- George Orwell, \emph{Nineteen Eighty-Four}
\end{flushright}

\noindent The core negative result is an \emph{escalating chain}: an \emph{a
fortiori} assumption on the attacker's knowledge, its decision-theoretic
consequence for a rational coercer, and the design bar the two force. Which links
are \emph{impossibilities} and which are \emph{incentive} claims is stated
throughout.

\subsection{Setup}
\label{sub:setup}

Fix the model of \S\ref{sec:model}. Write $K$ for the spend authority protecting
an asset worth $W$ to the holder $\Victim$, and $c$ for the expected cost to
$\Attacker$ of killing $\Victim$. Bitcoin makes the prize \emph{rivalrous} and
\emph{bearer}: the \emph{first} party to broadcast a valid spend takes $W$; one
prize, no appeal.

\begin{definition}[Seizable state, feasible path, racers]\label{def:racers}
The \textbf{seizable state} $\Sigma$ is everything $\Attacker$ obtains from the
encounter: seized device state, plus whatever it extracts from $\Victim$ during
the bounded coercion window. A \textbf{feasible path} from a state is an algorithm
outputting $K$ within the feasibility line (Stipulation~\ref{ass:feas}). If
$\Attacker$ releases $\Victim$, two \textbf{racers} may pursue $K$ --- $\Attacker$
from $\Sigma$, $\Victim$ from memory --- and $p_A \in [0,1]$ is the probability
$\Attacker$ spends \emph{first}. The prize being rivalrous, $\Victim$ spending
first to safety is $\Attacker$'s only loss condition once it holds \emph{any}
feasible path.
\end{definition}

\begin{definition}[Gray vs.\ clear outcomes]\label{def:gray}
An outcome is \textbf{gray} if afterward $0 < p_A < 1$: $\Attacker$ has a feasible
path yet $\Victim$, alive and free, can still win the race. The \textbf{clear}
outcomes are the endpoints, $p_A = 1$ and $p_A = 0$
(Figure~\ref{fig:deadly-race}).
\end{definition}

\begin{figure}[t]
\centering
\resizebox{\linewidth}{!}{%
\begin{tikzpicture}[node distance=7mm and 13mm]
  \node[gwevent] (seiz) {seize};
  \node[gwgate, right=of seiz] (g1) {feasible\\ path to $K$?};
  \node[gwend, above right=5mm and 13mm of g1] (fail) {clear\\ fail};
  \node[gwend, below right=5mm and 13mm of g1] (win)  {already\\ won};
  \node[gwtask, right=17mm of g1] (piv) {gray: $\Victim$\\ pivotal, $\pi>0$};
  \node[gwgate, right=of piv] (g2) {$W > c/\pi$?};
  \node[gwkill, right=of g2] (kill) {kill\\ $\Victim$};
  \node[gwend, below=9mm of kill] (surv) {$\Victim$\\ survives};
  \draw[gwflow] (seiz) -- (g1);
  \draw[gwflow] (g1) -- node[gwlabel, above left, inner sep=1pt]{$p_A{=}0$} (fail);
  \draw[gwflow] (g1) -- node[gwlabel, below left, inner sep=1pt]{$p_A{=}1$} (win);
  \draw[gwflow] (g1) -- node[gwlabel, above]{$0{<}p_A{<}1$} (piv);
  \draw[gwflow] (piv) -- (g2);
  \draw[gwflow] (g2) -- node[gwlabel, above]{yes} (kill);
  \draw[gwflow] (g2) -- node[gwlabel, right]{no} (surv);
\end{tikzpicture}}
\caption{\textbf{The Deadly Race decision.} At the \emph{clear} endpoints
($p_A\in\{0,1\}$) no race exists. In the \emph{gray} interval the holder is
\emph{pivotal} ($\pi>0$, Theorem~\ref{thm:safety}), so once $W>c/\pi$ a rational
$\Attacker$ has a positive incentive to \emph{assassinate} (Lemma~\ref{lem:drl}).
Clearing the Criterion means removing the gray branch entirely.}
\label{fig:deadly-race}
\end{figure}

The \textbf{priced} adversary $\mathrm{WDY}[W,c]$ is a WDY adversary additionally
endowed with $(W,c)$ and modeled as a cold expected-utility maximizer over
release, continued coercion and assassination; plain $\mathrm{WDY}$ carries no
such payoff, so the priced form is invoked only where a \emph{quantitative}
incentive is computed. Plain $\mathrm{WDY}$ worst-cases only the \emph{payoff}
coordinates ($W\!\to\!\infty$, $c\!\to\!0$), leaving \emph{capability} pinned at
the feasibility line and the registry-calibrated coercion window: sending either
to $\infty$ would collapse the floor or walk any time-gate to completion, a
category error rather than a stronger guarantee (Appendix~\ref{app:chain}).

\subsection{The Non-Negative-Proof Principle}
\label{sub:nnpl}

\begin{assumption}[Non-Negative-Proof Principle, NNPP]\label{lem:nnpl}
Kerckhoffs (Assumption~\ref{ass:kerckhoffs}) grants $\Attacker$ the
\emph{protocol} --- and that $\Victim$ \emph{may} hold brain-memory or an
undisclosed backup short-cutting $K$ --- but not how many such backups $\Victim$
keeps, or where. NNPP covers exactly that gap: what $\Attacker$ \emph{cannot} know
is whether any undisclosed backup \emph{actually} exists. The two do not conflict
--- one is the public mechanism, the other $\Victim$'s private state --- so NNPP
\emph{qualifies} Kerckhoffs, and we resolve the gap conservatively:
\textbf{$\Attacker$ always assumes an undisclosed backup exists}. Something as
small as an SD card, as inconspicuous as steganographic notes in a book, or as
interior as memorized email credentials suffices to hold a short-cutting state.
\end{assumption}

\noindent Both halves point the \emph{same} way --- toward killing --- since
each \emph{raises} $\Attacker$'s assessed payoff and hence its proclivity to kill;
each is therefore \emph{a fortiori}. In practice NNPP reduces to a plain
observation: \textbf{exhaustive search is so ineffective that $\Attacker$ always
assumes another backup exists, whatever the terrorized $\Victim$ pleads.} It is
rational, not paranoid --- non-possession is a \emph{negative existential} that
cannot be certified on either reading of ``prove,'' the search reading failing
because its last, irreducible site is the ``few cubic centimetres'' of the
epigraph. An $\Attacker$ that can read that space is a degenerate fixed point of
the model: under an uncontested overlord there is no custody left to resist, so
the only world falsifying NNPP is one in which the question does not arise
(Remark~\ref{rem:selfdefeat}).

\begin{corollary}[Cooperation cannot dispel the incentive]\label{cor:disclosure}
Let $\Victim$ hold $b$ resumable copies and, under the coercion oracle, disclose
some. Each disclosure \emph{confirms} a copy existed, raising a lower bound
$\hat b \le b$; none yields the upper bound $\hat b \ge b$, which would need the
covering search just ruled out. Disclosure is a \emph{monotone lower-bound
estimator} --- ``$\Victim$ has revealed everything'' is exactly the direction NNPP
forbids. Hence no degree of cooperation drives $\Attacker$'s posterior on
``$\ge 1$ undisclosed copy'' to zero: in any design leaving a released $\Victim$ a
resumable path, the incentive of \S\ref{sub:drl} cannot be discharged by
surrender, only by the victim's elimination.
\end{corollary}

\subsection{The Deadly-Race Lemma}
\label{sub:drl}

\begin{lemma}[Deadly-Race Lemma]\label{lem:drl}
Let a custody scheme admit a \emph{gray} coercion outcome
(Definition~\ref{def:gray}): on release, $\Attacker$ holds a feasible path to $K$
($p_A > 0$) while $\Victim$ retains one too ($p_A < 1$). Then
$\mathrm{WDY}[W,c]$ has a \textbf{strictly positive incentive to eliminate}
$\Victim$ (and any co-racing heirs) whenever $(1 - p_A)\,W > c$.
\end{lemma}

\begin{proof}
With $\Victim$ alive and free, $\Attacker$'s expected take from the race is
$p_A\,W$. By NNPP (Assumption~\ref{lem:nnpl}) $\Attacker$ cannot obtain or demand
a verified assurance that $\Victim$ will not race --- no such proof exists --- so
its \emph{only} instrument for pushing the win probability toward $1$ is to remove
the competing racer. Eliminating $\Victim$ drives that probability to
${\approx}\,1$ and the expected take to ${\approx}\,W$. The marginal gain is
$(1 - p_A)\,W$, and a rational $\Attacker$ acts on it exactly when it exceeds $c$.
In any gray outcome $p_A < 1$, so the gain is positive; for assets large relative
to $c$ the inequality holds.
\end{proof}

\noindent The result is \emph{decision-theoretic, not cryptographic}: the danger
is created by the \emph{shape} of the recovery path, not its hardness. And it
\emph{peaks where the marketing promises safety} --- a ``graceful'' fallback (a
slow-but-feasible recovery on a seized device, a time-locked rescue, a delayed
vault) holds $p_A$ well below $1$ while leaving $W$ fully at stake, i.e.\ it
\emph{maximizes} $(1-p_A)\,W$. Graceful degradation, a virtue everywhere else, is
here a liability --- hence the title.

\begin{remark}[Release is the trigger]\label{rem:release-trigger}
The incentive is a property of the released gray state: it is
$\Victim$-as-free-racer whom $\Attacker$ is paid to remove. A scheme that simply
hands $\Attacker$ the funds outright kills no one; the gray fallback is
\emph{worse than} that clean loss. And by Corollary~\ref{cor:disclosure} no amount
of surrender removes the incentive, so only the victim's elimination closes it.
\end{remark}

\subsection{The No-Gray-Area Criterion}
\label{sub:nga}

\begin{criterion}[No-Gray-Area Criterion]\label{prin:nga}
Under a fixed threat model, a coercion-\emph{safe} scheme must admit \textbf{no
gray outcome}: every encounter must resolve to a \emph{clear} endpoint ---
$p_A = 0$ (the attack plainly \emph{fails}: $\Sigma$ feasibly reaches nothing and
$\Victim$ holds nothing transmissible left to surrender) or $p_A = 1$ (it plainly
\emph{succeeds}: $\Attacker$ obtains $K$ within the encounter, so no
post-release race exists). The forbidden interval
$0 < p_A < 1$ is exactly where Lemma~\ref{lem:drl}'s incentive lives.
\end{criterion}

\noindent The Criterion is a \textbf{bar, not a construction}, and its two
endpoints are not equally desirable --- $p_A = 1$ costs $\Victim$ the funds --- but
both are \emph{survivable}: coercion resistance is properly a bar on \textbf{the
holder's safety} first and the funds only after. Nor need it stay a stipulation:
cast the encounter as a one-move game and no-gray \emph{is} the safety condition.

Fix a scheme $\Sigma$ and a priced adversary. In the one-move game
$\mathbf{Coerce}_{\Sigma}$ (Definition~\ref{def:game}, Appendix~\ref{app:game})
$\Attacker$ mounts an encounter (\ref{d:subjugation}--\ref{d:noperceptual}),
reaching a state with a set $R$ of \emph{racers} --- the released $\Victim$ and any
remote delegate --- then chooses $\textsf{spare}$ or $\textsf{kill}$, the latter
removing $\Victim$ from $R$ at cost $c$, and the race resolves with payoff
$p_A(S_a)\,W - c\,\mathbf 1[a=\textsf{kill}]$. Here $p_A(S)$ is $\Attacker$'s
probability of prevailing against racer set $S$, monotone in $S$, and by
reachability $\Attacker$ can eliminate only racers it can \emph{reach}, never a
remote delegate. $\Sigma$ is \textbf{coercion-safe} if no priced adversary, at any
reachable state, strictly prefers $\textsf{kill}$.

\begin{theorem}[No-Gray-Area, as a theorem]\label{thm:safety}
Write $\pi := p_A(R\setminus\{\Victim\}) - p_A(R) \ge 0$ for $\Victim$'s
\emph{pivotality}. $\Sigma$ is coercion-safe against every $\mathrm{WDY}[W,c]$
\emph{if and only if} $\pi = 0$ at every reachable state.
\end{theorem}

\begin{proof}
The kill increment is $\Delta U = \pi\,W - c$. ($\Leftarrow$) If $\pi = 0$
everywhere then $\Delta U = -c < 0$ and no priced adversary prefers
$\textsf{kill}$. ($\Rightarrow$, contrapositive) If $\pi > 0$ at some reachable
state then $\Delta U > 0$ for any $(W,c)$ with $W > c/\pi$; that adversary exists
and strictly prefers $\textsf{kill}$.
\end{proof}

\noindent An experiment-and-advantage restatement, with the clauses
\ref{d:kerckhoffs}--\ref{d:nodelegatecoerce} as oracles and \emph{reaching a gray
state} as the winning event, is in Appendix~\ref{app:game}. Pivotality has exactly
the two clearance routes of Corollary~\ref{cor:reach}: under \ref{r:removerace}
($R = \{\Victim\}$) removing $\Victim$ leaves no racer, so
$\pi = \mathbf 1[p_A>0] - p_A$, vanishing \emph{iff} $p_A \in \{0,1\}$ ---
Definition~\ref{def:gray}'s no-gray condition as a special case; under
\ref{r:remotewinner} killing $\Victim$ leaves the delegate $D$, so $\pi = 0$
\emph{iff $D$ dominates $\Victim$ as a racer}, and the state may be \emph{gray}
yet safe because the pivotal racer is unreachable. That is contingent --- a
co-located or merely \emph{slower} $D$ restores $\pi > 0$, which is the BitVault
case (\S\ref{sec:classify}). What the single-shot model leaves open is the
\emph{utility} behind $\Delta U$ (\S\ref{sec:limits}), not the equivalence.

\paragraph*{Grading protocols.}
Coercion-safety grades a protocol by how much of the adversary it withstands, and
\S\ref{sec:classify} uses that vocabulary: \textbf{naively} coercion-resistant if
it resists only a \emph{sub-Kerckhoffs} attacker, so that denial or decoy works
just while the attacker does not know $\Victim$'s setup; \textbf{weakly} if a full
WDY encounter does not by itself yield $K$ but \emph{may} leave a deadly-race
state ($\pi > 0$), the secret withheld \emph{during} coercion and the holder still
unsafe \emph{after} it; and \textbf{strongly} if it yields neither, which is
exactly coercion-safety. Cutting across the three, a protocol is
\textbf{time-gating} if reaching $K$ requires completing a gate of cost $G$ set
against the coercion budget $t$ --- a \emph{mechanism}, not a strength
(Definitions~\ref{def:naive-cr}--\ref{def:timegate}, Appendix~\ref{app:chain}).

\subsection{Reachability --- the two clearance routes}
\label{sub:reach}

\begin{corollary}[Reachability]\label{cor:reach}
Lemma~\ref{lem:drl}'s incentive is realizable only against a racer $\Attacker$
can \emph{reach}. A scheme can therefore satisfy Criterion~\ref{prin:nga} in
exactly two ways:
\begin{enumerate}[nosep,leftmargin=2.6em,label=(R\arabic*),ref=R\arabic*,font=\normalfont\bfseries\color{accentdark}]
  \item\label{r:removerace} \emph{remove the race} --- arrange that no feasible seizure path
        exists, forcing $p_A = 0$ (custody can stay \emph{individual}); or
  \item\label{r:remotewinner} \emph{remove $\Attacker$'s reach over the winner} --- make the party
        who \emph{will} win the race a \emph{remote delegate} beyond $\Attacker$'s
        coercive reach (custody becomes \emph{shared}, at the cost of a trusted
        party).
\end{enumerate}
\end{corollary}

\noindent These are the classification axis of \S\ref{sec:classify} ---
\emph{leave no race to run} vs.\ \emph{put the racer out of reach}. NNPP does
\emph{not} weaken the latter, as one might expect from $\Victim$'s inability to
prove no local coercible copy of the delegate's capability exists: that route's
holder-safety is exactly $\pi = 0$, and a hidden backup changes nothing, because
\textbf{assassination removes a racer but never extracts a secret}. What survives
is a difference about \emph{funds}, and a plainer objection --- an
access-dominating delegate \emph{is not self-custody}, its clearance resting on
$D$ refusing \textbf{a legitimate owner pleading for their life}
(\ref{d:nodelegatecoerce}), the same refusal that defeats $\Attacker$'s coerced
plea (Remark~\ref{rem:r2-cost}).

\subsection{Geographic, delegated, or tacit --- what forces TKBA}
\label{sub:fourprops}

Ask \emph{how} a time-gating protocol can attain $\pi = 0$ against a single WDY
encounter, confined to one \emph{site} (\ref{d:seizure}). There are exactly three
mechanisms, and each spends a different one of the \emph{Four Properties}:
\begin{enumerate}[nosep,leftmargin=2.4em,label=\textbf{(\arabic*)},ref=\arabic*]
  \item\label{p:know} knowledge-based authentication --- the secret in memory
        alone, i.e.\ \emph{non-geographic};
  \item\label{p:indiv} individual custody;
  \item\label{p:nonob} non-obscurity --- which \emph{is} the Kerckhoffs premise
        (\ref{d:kerckhoffs});
  \item\label{p:coerc} coercion resistance.
\end{enumerate}
A delegate is \textbf{access-dominating} (Definition~\ref{def:dominating}) if
its latency to reach and act on $K$ never exceeds $\Victim$'s, so removing
$\Victim$ leaves the race unchanged --- \ref{r:remotewinner}, at the cost of
(\ref{p:indiv}). A protocol is \textbf{geographic} (Definition~\ref{def:geographic})
if $K$'s gating is partitioned across sites, so a single-site encounter has no
feasible path --- custody may stay individual, but $K$ now requires locations, the
negation of (\ref{p:know}). It is \textbf{tacitly gated}
(Definition~\ref{def:tacit}) if it gates $K$ behind a \emph{non-transmissible}
in-head input (Assumption~\ref{ass:tkba}), which a single encounter can neither
extract nor reconstruct offline: $\pi = 0$ \emph{at a single site, with the secret
in memory only}, keeping both (\ref{p:know}) and (\ref{p:indiv}). The geographic
route clears the bar honestly but buys safety in \emph{physical} coin --- sites to
burgle, a victim to be moved between them --- so security rests on the physical
defenses: Bitcoin reduced to a \emph{glorified gold} guarded in vaults. The tacit
route escapes that regress. Appendix~\ref{app:chain} develops the geographic
attack surface, the \emph{literal} and \emph{cyber-physical} races, checkable
dominance, and why a perceptual oracle is worth its price only on the tacit route
(Remark~\ref{rem:po-variants}).

\begin{theorem}[DRL-safety trichotomy]\label{thm:trichotomy}
A time-gating custody protocol is DRL-safe against a WDY adversary only if, at
every reachable state, it is \textbf{access-dominating-delegated}
(\ref{r:remotewinner}), \textbf{geographic} (\ref{r:removerace}), or
\textbf{tacitly gated} (\ref{r:removerace}).
\end{theorem}

\noindent The proof (Appendix~\ref{app:chain}) turns on a two-blind-spots
enumeration: a co-located WDY that seizes every device and coerces the present
owner can be denied a piece of $K$'s gating only if that piece is \emph{elsewhere}
(geographic) or \emph{in the unextractable mind} behind a time-gate exceeding the
window (tacit); otherwise a racer at least as fast as $\Victim$ remains, which is
the delegate. Kerckhoffs having thrown obscurity out, these are the only ways to
restrict $K$'s availability at all. Of the three, the delegate forfeits
(\ref{p:indiv}) and the geographic route forfeits (\ref{p:know});
\textbf{only the tacit leg keeps all four Properties}, and TGPO (\S\ref{sec:gw})
is that construction. The equivalence
$\text{TKBA} \Leftrightarrow (\ref{p:know})\wedge(\ref{p:indiv})\wedge
(\ref{p:nonob})\wedge(\ref{p:coerc})$ is Corollary~\ref{cor:fourprops}
(Appendix~\ref{app:chain}): (\ref{p:know}) narrows the store to portable ones and
(\ref{p:coerc}) prunes those to the brain --- every carried device is seized at
the attack site, and a biometric is coerced --- then forces its content to be
non-transmissible. Unlike the trichotomy, it needs no exhaustiveness argument.

\subsection{A fifth property --- no hostage token}
\label{sub:hostage}

Devicelessness (\ref{p:know}) yields a fifth property the Four leave implicit:
immunity to indirect coercion \emph{through a seizable token}. A \emph{necessary}
token stored outside the holder's head is a \emph{hostage} --- even where seizing
$K$ is infeasible, $\Attacker$ can seize it, or threaten its destruction, and
trade it for the formal, irrevocable transfer of another good. The property is
orthogonal to \ref{r:remotewinner} but \emph{denied by construction} to any
geographic protocol, which needs precisely such tokens; only the in-head secret
has nothing to take hostage --- and, under Kerckhoffs, its holder's denial of a
backup is \emph{consistent with a perfectly ordinary setup} rather than a
suspicious claim to be tortured out of
(Remarks~\ref{rem:hostage}--\ref{rem:plausible-denial}). The qualifier is
load-bearing: leverage over a \emph{person} is a different mechanism, and no
custody design makes a relative unseizable (\S\ref{sub:empirical}).
